% Sample LaTeX file for PMLR double-column format
% This format is used for conferences that prefer double-column format.

\documentclass[pmlr,twocolumn]{jmlr} % This is the PMLR double-column format

% Recommended packages
\usepackage{booktabs} % For better tables
\usepackage{graphicx} % For including figures
\usepackage{url}      % For URLs
\usepackage{amsmath}  % For mathematical notation
\usepackage{algorithm2e} % For algorithms (optional)

% Proceedings information
\jmlrvolume{151}
\jmlryear{2024}
\jmlrworkshop{Artificial Intelligence and Statistics}

% Title and author information
\title{Sample Paper Title: A Comprehensive Approach to Statistical Learning}

% Use \Name{Author Name} to specify the name.
% If the surname contains spaces, enclose the surname
% in braces, e.g. \Name{John {Smith Jones}} similarly
% if the name has a "von" part, e.g \Name{Jane {de Winter}}.
% If the first letter in the forenames is a diacritic
% enclose the diacritic in braces, e.g. \Name{{\'E}louise Smith}

% Authors with diverse names demonstrating proper formatting:
\author{\Name{{\'E}lise {de la Fontaine}} \Email{elise.delafontaine@example.fr}\\
  \addr INRIA Paris\\
  Paris, France
  \AND
  \Name{Amit Kumar Singh} \Email{amit.singh@example.in}\\
  \addr Indian Institute of Technology\\
  Mumbai, India
  \AND
  \Name{Fatima Al-Rahman} \Email{fatima.alrahman@example.sa}\\
  \addr King Abdullah University\\
  Thuwal, Saudi Arabia
  \AND
  \Name{Yuki Tanaka} \Email{yuki.tanaka@example.jp}\\
  \addr University of Tokyo\\
  Tokyo, Japan
}

\editor{David Williams and Priya Sharma}

\begin{document}

\maketitle

\begin{abstract}
This is a sample abstract for a PMLR paper in double-column format. 
The abstract should be a brief summary of your paper, typically 
100-200 words. It should clearly state the problem, your approach, 
and the main results. Mathematical notation can be included, such as 
$f(x) = \sum_{i=1}^{n} w_i x_i$, and the abstract can span multiple 
lines as needed.
\end{abstract}

\begin{keywords}
Machine Learning, Sample Paper, PMLR Format, Two Column
\end{keywords}

\section{Introduction}

This is a sample PMLR paper in double-column format. The double-column
format is used for some PMLR proceedings, such as AISTATS. This format
uses the \texttt{jmlr} class with both the \texttt{pmlr} and 
\texttt{twocolumn} options, which automatically sets the header to 
read ``Proceedings of Machine Learning Research'' and formats the 
content in two columns.

You can include citations using standard \LaTeX{} citation commands,
such as \citet{example2023} or \citep{example2023}. Mathematical equations can be displayed:

\begin{equation}
  \label{eq:example}
  f(x) = \int_{-\infty}^{\infty} e^{-x^2} dx = \sqrt{\pi}
\end{equation}

In double-column format, equations and text automatically flow
within the column width.

\section{Methods}

Here you would describe your methods. You can include figures using
the standard \texttt{figure} environment. In two-column format, you
can use \texttt{figure} for single-column figures or \texttt{figure*}
for figures that span both columns:

\begin{figure}[htbp]
  \centering
  % Uncomment the line below when you have an actual figure
  % \includegraphics[width=\columnwidth]{figure.pdf}
  \caption{This is a single-column figure that fits within one column.
    Figures should be referenced in the text using 
    \texttt{Figure~\ref{fig:sample}}.}
  \label{fig:sample}
\end{figure}

For figures that need more space:

\begin{figure*}[htbp]
  \centering
  % Uncomment the line below when you have an actual figure
  % \includegraphics[width=0.8\textwidth]{wide-figure.pdf}
  \caption{This is a double-column figure that spans both columns.
    Use \texttt{figure*} for wide figures. Note that these figures
    are typically placed at the top of a page.}
  \label{fig:wide}
\end{figure*}

\subsection{Algorithms}

You can also include algorithms. Here's a simple example:

\begin{algorithm}[htbp]
  \SetAlgoLined
  \KwIn{Data $(x_1, \ldots, x_n)$}
  \KwOut{Parameters $\theta$}
  Initialize $\theta$ randomly\;
  \While{not converged}{
    $\nabla L(\theta) \leftarrow$ gradient\;
    $\theta \leftarrow \theta - \eta \nabla L$\;
  }
  \Return $\theta$\;
  \caption{Sample stochastic gradient descent algorithm}
  \label{alg:sample}
\end{algorithm}

\section{Results}

Present your results here. Tables can be formatted using the
\texttt{booktabs} package. Similar to figures, you can use
\texttt{table} for single-column tables or \texttt{table*}
for tables that span both columns:

\begin{table}[htbp]
  \centering
  \caption{Sample results (single column)}
  \label{tab:results}
  \begin{tabular}{lc}
    \toprule
    Method & Accuracy \\
    \midrule
    Baseline & 85.2\% \\
    Our Method & 92.1\% \\
    \bottomrule
  \end{tabular}
\end{table}

\begin{table*}[htbp]
  \centering
  \caption{Detailed results spanning both columns}
  \label{tab:detailed}
  \begin{tabular}{lccccc}
    \toprule
    Method & Accuracy & Precision & Recall & F1-Score & Training Time (s) \\
    \midrule
    Baseline & 85.2\% & 83.1\% & 87.3\% & 85.1\% & 120 \\
    Method A & 89.5\% & 88.2\% & 90.8\% & 89.5\% & 110 \\
    Our Method & 92.1\% & 91.5\% & 92.7\% & 92.1\% & 95 \\
    \bottomrule
  \end{tabular}
\end{table*}

\section{Conclusion}

Summarize your findings and their implications here. The double-column
format allows for more compact presentation while maintaining
readability.

\section*{Acknowledgments}

Acknowledge funding sources and collaborators here. Note that this
section is unnumbered (using \texttt{section*}).

% Bibliography
% PMLR uses author-year citations by default
% Use \citet{} for textual citations: "Smith and Jones (2023) showed..."
% Use \citep{} for parenthetical citations: "as shown previously (Smith and Jones, 2023)"

% You can use BibTeX (recommended) or write entries manually as shown below
% For BibTeX, use: \bibliography{yourbibfile}

\begin{thebibliography}{10}

\bibitem[Author and Collaborator(2023)]{example2023}
Author, A. and Collaborator, B.
\newblock Sample Paper Title.
\newblock \emph{Proceedings of Machine Learning Research}, 123:456--789, 2023.

\end{thebibliography}

% Appendix (optional)
\appendix
\section{Additional Details}

Any supplementary material, proofs, or additional details can be
included in an appendix. In two-column format, appendix content
also flows across both columns.

\end{document}

